\chapter*{Acknowledgements}

The six years of my PhD have been a long, winding, and unforgettable journey.
Looking back on this period, I know deeply that this thesis was not completed by me alone.
It grew from the help, support, companionship, patience, and kindness of many people.
I would like to express my sincere gratitude to everyone who has helped me along the way.

This research was supported by the Commonwealth through an Australian Government Research Training Program (RTP) Scholarship.

First, I would like to thank my supervisor, Professor Catherine Stampfl.
She is one of the most inclusive and supportive people I have ever met.
She provided me with comprehensive and thoughtful guidance,
and introduced me to the profound and fascinating field of condensed matter physics.
My foundation was not strong, yet I often aimed too high and could easily lose myself in overly abstract questions.
Cathy, as we all call her, always brought these questions back to concrete physical pictures and feasible research paths.
She was enthusiastic in guiding me and always treated each of my questions seriously.
In every discussion, I could feel her patience, rigour, and kindness.
For this, I remain deeply grateful.

I would also like to give special thanks to Dr Oliver James Conquest.
As a senior colleague in my condensed matter physics research,
he provided patient and continuous support in technical matters, computational details, and research practice.
His extensive technical help and detailed guidance were essential for me to complete my PhD research. 
More importantly, during my most difficult period,
he also offered practical help and sincere support.
These acts of kindness meant a great deal to me.

I also thank Ruoning Ji. She accompanied me through a difficult period and helped me recover courage that I had once lost.
She was like someone who suddenly appears in the storyline and pulls the protagonist back from a side quest to the main path. Without her companionship,
many days would have been much heavier. For this, I am very grateful.

During my PhD, many people gave me important help.
I would like to thank Dr Carla Verdi.
Although she later moved to the University of Queensland and our time together was limited, 
her knowledge, working style, and attitude toward life still had a deep influence on me.
I also thank the other members of the USyd CMT group, including but not limited to Dr Marco Fronzi, Dr Hongyang Ma, Dr Giyeok Lee, Naiyuan Dong, Keri Liang, and Jiarong Li.
Thank you for your companionship, discussions, and help during this important period.

In addition, I would like to thank my GP, Dr Joo Faa Yee.
Over the past six years, he has given me tremendous help, understanding, and support.
This support accompanied me through many difficult moments during my PhD.

I would also like to thank my friends in real life, including but not limited to Rui Yu and Jiawei Sun.
Sydney is a city with almost no winter, and your presence made this city even warmer.
Thank you for your companionship, encouragement, and understanding at different stages.
Meanwhile, I also thank the friends I met online.
Dr Yulin Gong from the School of Mathematics at the University of Bristol patiently gave me guidance and helped me become less restless on many questions, from which I benefited greatly.
Jiarong Li from the Department of Physics at Tsinghua University and Weiyan Gao from the School of Materials and Chemical Technology at Institute of Science Tokyo are my best friends in cyberspace.
Although we have not yet met in person, we can talk about almost anything.
Such friendship is equally precious.

I would like to thank my family, especially my parents. Thank you for your constant support, understanding, and tolerance over the past six years.
Many choices during my PhD were not easy, but they always supported me from behind and allowed me to keep moving forward.

I would also like to leave my thanks to many presences that do not usually belong in acknowledgements.
The kookaburras that screamed loudly outside my window, the lovely cockatoo that once bit through my fingernail, 
the Galah that appeared by chance and flashed past like a little pink fairy,
and the always somewhat silly-looking pelicans all coloured my life in their own ways.
The jacarandas laid down a violet sea of flowers between spring and summer,
the Gymea lilies stood tall and graceful,
the Blue Mountains opened into waterfalls and morning mist,
the Barrenjoey Lighthouse on the hill at Palm Beach looked out alone toward the sea,
the Kiama blowhole surged beneath the rainbow,
and the towering Norfolk Island pines stood quietly by the coast.
The brilliant southern night skies and aurora,
together with Anytime Fitness Newtown, where I went most often during these six years,
also witnessed this winding, lonely, and unforgettable period.
In quiet or lively ways, they accompanied me, supported me, and reminded me that life is always broader than a difficult period.

Finally, I would also like to thank the person whom I may not yet know,
but whom I am still willing to believe I may meet somewhere in the future.
The real world may not have a closed-form analytical solution, and many paths cannot be known in advance.
One thing these six years of PhD study have taught me is that every journey is a path integral.
Perfection does not exist, and the road ahead cannot yet be fully known.
Precisely because of this, this winding period has become one of the most unforgettable parts of my life.
